\chapter{Constitutive modelling}
\label{ch:modelling}%
% Aim: 15-20 pages. From the measurements to a law that can be used in a numerical analysis.

\section{Modelling framework}
\label{sec:framework}
% Assumptions: oedometric strain path, isothermal tests, state described by porosity and ice fraction.
% General form of the densification law:
% \dot{\varepsilon} = A \, \sigma^{n} \exp(-Q/RT) \, f(\phi, \text{ice fraction})

\section{Calibration}
\label{sec:calibration}
\subsection{Stress dependence}
\subsection{Temperature dependence}
\subsection{Dependence on porosity}
\subsection{Dependence on ice content}
\subsection{Parameter uncertainty}
% Propagate the uncertainty on seal friction into the stress exponent and the pre-factor.

\section{Comparison with existing laws}
\label{sec:comparison}
% Glen-type flow law for ice, firn densification laws, frozen-soil creep models.

\section{Formulation for numerical implementation}
\label{sec:implementation}
% From the uniaxial law to a three-dimensional formulation: additional assumptions required,
% temperature as an input field, elastic part, integration scheme.
\subsection{Verification}
% Single-element simulation of an oedometric test compared with the experiments.

\section{Limitations and range of validity}
\label{sec:model_limits}
% Radial stress not measured; no true steady state; extrapolation to lunar cryogenic temperatures;
% one simulant and one ice morphology.

\section{Implications for lunar applications}
\label{sec:implications}
% Order-of-magnitude consequences for compaction under foundations, excavation and long-term stability.
